

The intersection of large language models, multi-agent systems, and blockchain oracle design has seen rapid development, but no prior work has systematically evaluated multi-LLM oracle architectures for prediction market resolution. This section surveys the most relevant research threads and identifies gaps in existing literature.

\subsection{LLM-Based Oracle Systems}

The most directly relevant prior work is Chainlink Labs' 2025 study evaluating a single-LLM oracle built with the DSPy framework~\cite{chainlink2025aioracles}. Their system uses GPT-4o across three modules: a Question Transformer employing chain-of-thought reasoning, an Information Gathering module performing web scraping with full webpage content retrieval, and a Fact Extraction module producing explicit reasoning traces. Evaluated on 1,660 resolved Polymarket outcomes each with over \$100,000 in trading volume, the system achieved 89.3\% overall accuracy, with dramatic category-dependent variation: 99.7\% on sports, 85.0\% on cryptocurrency, and 84.3\% on political outcomes. The primary bottleneck identified was \emph{temporal cognition}: LLMs struggle with relative time expressions and maintaining awareness of when information was created relative to events being verified. Sports questions, with their discrete temporal endpoints, dramatically outperform politics and cryptocurrency, where temporal boundaries are fuzzy and developments are ongoing. Notably, the study did not filter for potential training data contamination, leaving open the possibility that some accuracy reflects memorized rather than retrieved information. Additionally, their results are also possible skewed given the fact that they only resolved markets with high trading volumes, meaning there are a plethora of information about these markets

UMA's OOTruthBot represents a deployed AI oracle integration that evolved from a single Perplexity API call to a three-layer architecture~\cite{uma2025truthbot}. A lightweight Router dispatches questions to specialized solvers: a Code-Runner for deterministic API queries and a Perplexity-based solver for open-ended questions. An Overseer cross-references outputs against Polymarket live prices as a quality filter. The system achieves approximately 78\% overall accuracy (up to 95\% on objective events with clear data) at roughly \$0.005 per resolution, orders of magnitude cheaper than human-mediated resolution. A notable safety feature is the ``Too Early'' output when evidence is insufficient, implementing a principled abstention mechanism that prevents the system from guessing when it lacks confidence.

\subsection{Multi-Agent Architectures for Improved Reasoning}

The use of multiple LLM instances to improve accuracy has been explored through
several architectural paradigms that differ in how agents interact: independently
(ensemble voting), iteratively (debate and consensus), or hierarchically (layered
aggregation).

The foundational work on multi-agent debate by Du et al.\ demonstrated that
structured argumentation between LLM instances improves factual accuracy and
reasoning~\cite{du2023debate}. Using three agents over two rounds of debate, they
observed arithmetic accuracy improvements from approximately 70\% to 95\%. Two
findings are particularly relevant to oracle design. First, providing full reasoning
chains during debate substantially outperforms sharing only final answers, as agents
benefit from understanding \emph{why} other agents reached their conclusions. Second,
cross-model debate is effective: ChatGPT and Bard, each producing incorrect
individual answers, reached correct solutions through debate. Both observations
motivate our deliberative consensus architecture, which shares full reasoning traces
across diverse model families.

Omar et al.\ introduced the Iterative Consensus Ensemble (ICE), where three LLMs
critique each other iteratively until convergence~\cite{omar2025ice}. ICE improved
overall accuracy from 60.2\% to 74.0\% across over 4,000 questions, with most
questions settling within 2--3 rounds. This convergence behavior suggests that
extended deliberation yields diminishing returns, an observation relevant to the
cost-accuracy tradeoffs central to oracle system design.


\subsection{When Multi-Agent Approaches Fail}
The most important counterpoints to multi-agent enthusiasm come from recent empirical investigations of failure modes. Recent work formalized a ``tyranny of the majority'' effect: as the number of agents providing the same answer increases, whether correct or not, minority agents conform, creating echo chambers rather than enabling error correction~\cite{estornell2024tyranny}. Three interventions were proposed: diversity-pruning, quality-pruning, and misconception-refutation. Complementary work found that group accuracy frequently degrades over debate rounds, and that even introducing a weaker model into a majority of stronger models can diminish overall performance~\cite{talkisntcheap2025}. A large-scale trace analysis examined over 1,600 annotated traces across seven multi-agent frameworks, identifying 14 unique failure modes in three categories, with inter-agent misalignment as the dominant failure type~\cite{cemri2025multiagent}. Using 16$\times$ more compute reduced failure rates by only 43\%, while hierarchical structures outperformed flat architectures by 10--15\%.

The connection to group decision-making theory is instructive. A foundational distinction separates ``demonstrable'' tasks, where a correct answer can be logically proven, from ``judgmental'' tasks where no objectively verifiable solution exists~\cite{laughlin2011group}. For demonstrable tasks, a ``Truth Wins'' dynamic allows even a single correct minority member to persuade the group; for judgmental tasks, a ``Majority Wins'' dynamic dominates regardless of correctness. Oracle resolution straddles both categories: sports outcomes with official results resemble demonstrable tasks, while politically ambiguous questions resemble judgmental ones where convergence may reflect shared biases rather than accuracy.

\subsection{Correlated Errors and the Limits of Aggregation}

Ensemble methods rely on error independence. The Condorcet Jury Theorem guarantees that majority voting improves accuracy only when individual errors are uncorrelated. However, Kim et al.\ find that LLMs frequently agree on the same incorrect answer, indicating substantial error correlation across models~\cite{kim2025correlated}. Related work documents homogenization effects introduced by alignment training~\cite{wu2024monoculture, bai2025systematic}.

These results imply that naive majority voting may yield limited robustness gains. A key question is whether structured deliberation can reduce correlated error or whether it amplifies shared blind spots.

\subsection{Hybrid AI-Human Systems and Principled Escalation}

The question of when an AI system should defer to human judgment is formalized in the learning-to-defer literature, which provides the theoretical grounding for our escalation framework.

Mozannar and Sontag provided a consistent surrogate loss that jointly learns a classifier and rejector~\cite{mozannar2020consistent}. The Bayes-optimal deferral rule has an intuitive form: defer when the probability that the expert is correct exceeds the model's maximum posterior probability for any class. In oracle terms, this means deferring to human arbitration when the AI system's confidence is lower than expected human accuracy, a principled cost-benefit calculation rather than an ad-hoc threshold.

Bansal et al.\ found that AI explanations did not increase human-AI team accuracy beyond simply showing confidence scores, as explanations increased acceptance of AI recommendations regardless of correctness~\cite{bansal2021does}. This finding cautions against assuming that providing LLM reasoning traces to human arbitrators in an escalation pipeline will improve their judgment; confidence signals and disagreement patterns may be more actionable.

Selective prediction formalizes the tradeoff between accuracy and coverage: a system that declines to answer uncertain questions will be more accurate on those it does answer, but at the cost of resolving fewer questions automatically~\cite{chow1970optimum, geifman2017selective}. UMA's OOTruthBot already implements principled abstention through its ``Too Early'' output~\cite{uma2025truthbot}. Recent work formalizes three-tier LLM-human architectures where a base model routes to a large model which routes to a human expert~\cite{tiered2025architecture}, a pattern that maps directly to oracle system design. Our escalation framework builds on these approaches by using multi-agent disagreement and confidence as principled routing signals.


\subsection{Gaps in Existing Research}


Across these threads, no work has systematically evaluated multi-LLM
architectures for prediction market resolution, characterized error
correlation in oracle settings specifically, or developed principled
escalation criteria grounded in the learning-to-defer framework. The
methodology in Section~\ref{methodology} addresses these gaps through a controlled
comparison of independent aggregation and deliberative consensus, an
error correlation analysis across model pairs, and an escalation framework
that uses inter-agent disagreement and confidence as routing signals.